\subsection{实数的近似}

定义 1. 给定一个数 \(\varepsilon > 0\) ，实数 \(r\) 称为精确到 \(\varepsilon\) 的实数 \(x\) 的近似，如果 \(|x - r| \leqslant \varepsilon\) ；\(r\) 称为不足近似，如果 \(r \leqslant x\) ；称为过剩近似，如果 \(r \geqslant x\) 。

设 A 是 R 的一个稠密子集。对每个 \(x \in R\) 与每个 \(\varepsilon > 0\) ，存在 x 的属于 A 的精确到 \(\varepsilon\) 的不足（相应地，过剩）近似，因为区间 \(]x - \varepsilon, x[\) （相应地，\(]x, x + \varepsilon[\) ）至少含有 A 的一个点。现在若考虑一个给定的严格递减序列 ( \(\varepsilon_{n}\) ) （由数 \textgreater{} 0 组成，趋于 0），而 \(r_{n}\) 是 x 的精确到 \(\varepsilon_{n}\) 的近似，则序列 ( \(r_{n}\) ) 当 n 趋于无穷时以 x 为极限。

在 A 是加法群 R 的子群且限制 \(\varepsilon_{n}\) 属于 A 的情形，我们对每个 \(x \in R\) 都能按典型方式定义一个由 x 的属于 A 的不足近似组成的序列 \((r_{n})\) 。

因为由 Archimedes 公理（§ 2，第 1 节，定理 1），整数 \(p\) （满足 \(p\varepsilon_n \leqslant x\) ）组成的集合有一个最大元 \(p_n\) ；换言之，存在唯一的整数 \(p_n\) 使得

\[
p _ {n} \varepsilon_ {n} \leqslant x <   (p _ {n} + 1) \varepsilon_ {n}.\tag{1}
\]

由于 \(|x - p_n\varepsilon_n| \leqslant \varepsilon_n\) ，由此推出 \(p_n\varepsilon_n\) 是 \(x\) 的精确到 \(\varepsilon_n\) 的不足近似，且由假设属于 A；类似地，\((p_n + 1)\varepsilon_n\) 是 \(x\) 的属于 A 的精确到 \(\varepsilon_n\) 的过剩近似，而两个序列 \((p_n\varepsilon_n)\) 与 \(((p_n + 1)\varepsilon_n)\) 都以 \(x\) 为极限。

